\item Dos ondas simultáneamente presentes sobre una cuerda larga tienen una diferencia de fase $\phi$ entre ellas de modo que la onda formada a partir de su combinación está descrita por
$$ y(x, t) = 2A \sin\left(kx + \frac{\phi}{2}\right) \cos\left(\omega t - \frac{\phi}{2}\right) $$
\begin{enumerate}
    \item A pesar de la presencia del ángulo de fase $\phi$, ¿es cierto que los nodos están separados por media longitud de onda? Explique.
    \item ¿Son los nodos diferentes de alguna manera de lo que serían los nodos si $\phi$ fuera cero? Explique.
\end{enumerate}